\rok{Новембар 2023.}{H}

Први практични тест из Алгоритама и структура података

\zadatak{1}{Брисање елемената из стека (Stack Pop)}
Написати методу која имплементира стандардан начин за брисање елемената из стека.

\textbf{Додатне напомене за решавање задатка:}
\begin{itemize}
    \item Алгоритам написати у оквиру закоментарисане методе pop?().
    \item Поменута метода се налази у оквиру класе \texttt{Stack} која се налази у придруженом шаблону.
    \item Обезбедити код у \texttt{Main} методи за тестирање алгоритма.
\end{itemize}



\zadatak{2}{Да ли елементи двоструко спрегнуте листе чине Армстронгов број?}
Креирати функцију која ће проверавати да ли су елементи двоструко спрегнуте листе Армстронгов број. Армстронгов број је број који је збир властитих цифара од којих је свака подигнута на степен укупног броја цифара. Алгоритам је неопходно написати са линеарном временском комплексношћу $O(n)$.

\textbf{Додатни захтеви за решавање задатка:}
\begin{itemize}
    \item Захтеван потпис методе:
    \begin{Verbatim}[fontsize=\footnotesize, numbers=left, xleftmargin=10mm]
public static bool IsNumberArmstrong ()
    \end{Verbatim}
\end{itemize}

\textbf{Примери:}
\begin{itemize}
    \item \textbf{Улаз:} $1 \leftrightarrow 5 \leftrightarrow 3$ \\
          Објашњење: $1^3 + 5^3 + 3^3 = 1 + 125 + 27 = 153$ \\
          \textbf{Излаз:} \texttt{true}
    \item \textbf{Улаз:} $1 \leftrightarrow 3$ \\
          Објашњење: $1^2 + 3^2 = 1 + 9 = 10$ \\
          \textbf{Излаз:} \texttt{false}
\end{itemize}



\zadatak{3}{Исограм}
Креирати функцију која проверава да ли је неки број „исограм". Исограм је број који не дуплира цифре. Алгоритам је неопходно написати са линеарно-логаритамском временском комплексношћу $O(n \log n)$.

\textbf{Додатни захтеви за решавање задатка:}
\begin{itemize}
    \item Алгоритам треба да исписује захтевани текст.
    \item Као помоћ може се користити алгоритам за сортирање који је дат у шаблону задатка.
    \item Цифре прослеђеног броја неопходно је сместити у низ. Уколико се користи динамички низ, користити функцију \texttt{ToArray()} kako би се од њега добио статички низ.
    \item Осим динамичког низа, може се користити и структура једноструко повезане листе која је дата у оквиру задатка 2.
    \item Захтеван потпис методе:
    \begin{Verbatim}[fontsize=\footnotesize, numbers=left, xleftmargin=10mm]
public void IsNumberIsogram(long number)
    \end{Verbatim}
\end{itemize}

\textbf{Примери:}
\begin{itemize}
    \item \textbf{Улаз 1:} \texttt{12345} \\
          \textbf{Излаз 1:} \texttt{"Number 12345 is isogram!"}
    \item \textbf{Улаз 2:} \texttt{12324} \\
          \textbf{Излаз 2:} \texttt{"Number 12324 is NOT isogram!"}
\end{itemize}

\sablon{Шаблон првог теста новембар 2023. група H}{2023/Novembar/GrupaH/AISP2023_Test1_GrupaH.linq}
